% \documentclass[11pt,a4]{article}

% \usepackage[margin=2cm]{geometry}

% \usepackage{amsmath}
% \usepackage{url}
% \usepackage{amsfonts}
% \usepackage{amssymb}

% \usepackage[utf8]{inputenc}

% % Default fixed font does not support bold face
% \DeclareFixedFont{\ttb}{T1}{txtt}{bx}{n}{10} % for bold
% \DeclareFixedFont{\ttm}{T1}{txtt}{m}{n}{10}  % for normal

% % Custom colors
% \usepackage{color}
% \definecolor{deepblue}{rgb}{0,0,0.5}
% \definecolor{deepred}{rgb}{0.6,0,0}
% \definecolor{deepgreen}{rgb}{0,0.5,0}

% \usepackage{listings}

% % Python style for highlighting
% \newcommand\pythonstyle{\lstset{
% language=Python,
% basicstyle=\ttm,
% otherkeywords={self},             % Add keywords here
% keywordstyle=\ttb\color{deepblue},
% emph={MyClass,__init__},          % Custom highlighting
% emphstyle=\ttb\color{deepred},    % Custom highlighting style
% stringstyle=\color{deepgreen},
% frame=tb,                         % Any extra options here
% showstringspaces=false            %
% }}


% % Python environment
% \lstnewenvironment{python}[1][]
% {
% \pythonstyle
% \lstset{#1}
% }
% {}

% % Python for external files
% \newcommand\pythonexternal[2][]{{
% \pythonstyle
% \lstinputlisting[#1]{#2}}}

% % Python for inline
% \newcommand\pythoninline[1]{{\pythonstyle\lstinline!#1!}}


% \usepackage{collectbox}

% \newcommand{\mybox}[2]{$\quad$\fbox{
% \begin{minipage}{#1cm}
% \hfill\vspace{#2cm}
% \end{minipage}
% }}


% \usepackage{fancyhdr}
% \pagestyle{fancy}
% \rhead{Programmazione 1 - Esame del 24 giugno 2020}

% \include{book_header}

%\begin{document}
\thispagestyle{empty}
\hrule
\begin{center}
   {\Large {\bf Programmazione 1 \hspace{3cm} Esame del 24 giugno 2020}}
\end{center}
{\bf Cognome: }\hspace{2.5cm} {\bf Nome: } \hspace{2.5cm} {\bf Matricola: } \hspace{2.5cm} {\bf Postazione: } \\\
\hrule


%%%%%%%%%%%%%%%%%%%%%%%%%%%%%%%%%%%%%%%%%%%%%%%%%%%%%%%%%%%%%%%%%%%%%%%%%%%%%
\section*{Prima Parte}
Nella prima parte dell'esame dovete usare esclusivamente le funzioni ricorsive e le {\tt pairlists},
come studiate nella prima parte del corso. Dovete importare la libreria {\tt pairlists} che viene fornita nella vostra
cartella d'esame. Dovete svolgere i seguenti esercizi, completando il file {\tt soluzione\_template.py}
e {\underline {\bf utilizzando le funzioni di test che trovate nel blocco {\tt main} in coda al file}}.
\begin{enumerate}
%%%%%%%%%%%%%%%%%%%%%%%%%%%%%%%%%%%%%%%%%%%%%%%%%%%%%%%%%%%%%%%%%%%%%%%%%%%%%
\item {\bf (Punti 3)} Scrivere una funzione {\tt Norma(x, t)} che prende in input una lista di numeri {\tt x} (ovvero un vettore di $\mathbb{R}^n$) e un intero $t$ e restituisce la norma del vettore. Se $t=0$ restituisce la norma infinito, ovvero $\max_{i}|x_i|$; se $t=1$ calcola la norma 1, ovvero $\sum_i |x_i|$;
se $t=2$ la norma 2, ovvero $\sqrt{\sum_i x_i^2}$.

%%%%%%%%%%%%%%%%%%%%%%%%%%%%%%%%%%%%%%%%%%%%%%%%%%%%%%%%%%%%%%%%%%%%%%%%%%%%%
\item {\bf (Punti 4)} Scrivere una funzione {\tt Integrale(F, a, b, dx)} che calcoli l'integrale definito di una funzione $f$ tra i limiti $a$ e $b$, che può essere approssimato numericamente usando la formula:
    $$
    \int_a^b f(x)\, dx = [f(a+dx/2) + f(a+dx/2+dx) + f(a+ dx/2 + 2dx )+ ...]dx
    $$
per valori piccoli di $dx$.

%%%%%%%%%%%%%%%%%%%%%%%%%%%%%%%%%%%%%%%%%%%%%%%%%%%%%%%%%%%%%%%%%%%%%%%%%%%%%
\item {\bf (Punti 6)} Scrivere una funzione che implementa l'algoritmo di {\it merge-sort} utilizzando le {\tt pairslist}.
La funzione {\tt MergeSort(Ls, Compare)} da implementare, prende in input una lista di elementi (lista di tipo {\tt pairslist}),
e una predicato di confronto {\tt Compare(x, y)} tra due elementi della lista,
e restituisce in output una nuova lista ordinata in base alla funzione di confronto {\tt Compare(x, y)} data in input.
\end{enumerate}


%%%%%%%%%%%%%%%%%%%%%%%%%%%%%%%%%%%%%%%%%%%%%%%%%%%%%%%%%%%%%%%%%%%%%%%%%%%%%
\section*{Seconda Parte}
Nella seconda parte dell'esame potete usare qualsiasi libreria di Python.

Si consideri un file di testo, chiamato {\tt dpc-covid19-ita-regioni-latest.csv}, che contiene i
dati della protezione civile riguardo l'epidemia in corso, rilasciati il 21 giugno 2020 su GitHub.
Per ogni riga del file, tranne la prima che contiene la descrizione di ogni campo, vengono riportate le seguenti informazioni:
\begin{verbatim}
 0, data: data a cui si riferiscono le informazioni
 1, stato: codice dello stato
 2, codice_regione: codice numerico per la regione
 3, denominazione_regione: nome della regione
 4, lat: latitudine
 5, long: longitudine
 6, ricoverati_con_sintomi: numero di pazienti ricoverati con sintomi
 7, terapia_intensiva: numero di pazienti in terapia intensiva
 8, totale_ospedalizzati: numero di pazienti ospedalizzati
 9, isolamento_domiciliare:	numero di pazienti in isolamento domiciliare
10, totale_positivi: numero di pazienti positivi in totale dall'inizio dell'epidemia
11, variazione_totale_positivi:	numero di pazienti positivi rispetto al giorno precedente
12, nuovi_positivi: numero di nuovi pazienti positivi per la data di riferimento
13, dimessi_guariti: numero di pazienti guariti
14, deceduti: numero di pazienti deceduti
15, totale_casi: totale dei casi registrati
16, tamponi: totale dei tamponi eseguiti
17, casi_testati: numero totale dei casi testati
18, note_it: eventuali note in italiano
19, note_en: eventuali note in inglese
\end{verbatim}


\begin{enumerate}
\setcounter{enumi}{3}

%%%%%%%%%%%%%%%%%%%%%%%%%%%%%%%%%%%%%%%%%%%%%%%%%%%%%%%%%%%%%%%%%%%%%%%%%%%%%
\item {\bf (Punti 2)}
Si scriva una classe {\tt Region} con i seguenti attributi: {\tt Data}, {\tt Nome}, {\tt TerapiaIntensiva}, {\tt NuoviPositivi}, {\tt Deceduti}, {\tt Totale}, e  {\tt Tamponi}.
Tranne il nome che è una stringa, tutti gli altri attributi sono dei numeri interi.
La classe deve contenere almeno un metodo costruttore {\tt \_\_init\_\_(self, row)}, che data una riga letta dal fila,
inizializza un oggetto corrispondente di tipo {\tt Region}.

%%%%%%%%%%%%%%%%%%%%%%%%%%%%%%%%%%%%%%%%%%%%%%%%%%%%%%%%%%%%%%%%%%%%%%%%%%%%%
\item {\bf (Punti 1)} Si implementi il metodo {\tt \_\_str\_\_(self)}, che restituisce una stringa contenente i primi tre attributi, nel formato seguente:

{\tt "Regione: Lombardia, Totale: X, Deceduti: X, Tamponi: X"}


%%%%%%%%%%%%%%%%%%%%%%%%%%%%%%%%%%%%%%%%%%%%%%%%%%%%%%%%%%%%%%%%%%%%%%%%%%%%%
\item {\bf (Punti 2)} Si implementi il metodo {\tt \_\_lt\_\_(self, other)}, che restituisce True se la regione {\tt self} ha avuto
un rapporto tra il numero di morti e il totale dei pazienti infetti inferiore allo stesso rapporto dell'altra regione codificata
nell'oggetto {\tt other}.

%%%%%%%%%%%%%%%%%%%%%%%%%%%%%%%%%%%%%%%%%%%%%%%%%%%%%%%%%%%%%%%%%%%%%%%%%%%%%
\item {\bf (Punti 3)} Si scriva una funzione {\tt Parse(filename)}, che prende in input il nome di un file,
e restituisce una lista di oggetti di tipo {\tt Region}, un oggetto per ogni riga del file (a parte la prima riga di intestazione).


%%%%%%%%%%%%%%%%%%%%%%%%%%%%%%%%%%%%%%%%%%%%%%%%%%%%%%%%%%%%%%%%%%%%%%%%%%%%%
\item {\bf (Punti 3)} Si scriva una funzione {\tt TopRegion(Ls, n=5)}, che prende in input una lista di oggetti di tipo {\tt Region}
come restituiti dalla funzione {\tt Parse}, e che restituisce le prime $n$ regioni con il maggior rapporto tra numeri di morti e pazienti infetti.


%%%%%%%%%%%%%%%%%%%%%%%%%%%%%%%%%%%%%%%%%%%%%%%%%%%%%%%%%%%%%%%%%%%%%%%%%%%%%
\item {\bf (Punti 3)} Si scriva una funzione {\tt ParseFiles(names)} che prende in input una lista di nomi di files che devono
essere letti. La funzione per ogni nome della lista, deve richiamare la funzione {\tt Parse}, leggere il contenuto del file ottenendo
una lista di oggetti di tipo {\tt Region}, e deve restituire alla fine la lista complessiva di oggetti letti dai file di input.

Nella funzione di test, si chiede di leggere cinque file, dello stesso tipo del file {\tt .csv} descritto sopra,
ma che si riferiscono a date diverse.

%%%%%%%%%%%%%%%%%%%%%%%%%%%%%%%%%%%%%%%%%%%%%%%%%%%%%%%%%%%%%%%%%%%%%%%%%%%%%
\item {\bf (Punti 5)} Si scrive una funzione {\tt ComputeAverage(Ls)}, che prende in input una lista di oggetti di tipo {\tt Region},
letti per esempio con la funzione {\tt ParseFiles(names)}, e calcola, per ogni regione, la media aritmetica giornaliera dei nuovi casi registrati.
La funzione deve restituire un dizionario con una chiave per ogni regione, a cui corrisponde come valore la media giornaliera calcolata di nuovi casi.


\end{enumerate}

%%%%%%%%%%%%%%%%%%%%%%%%%%%%%%%%%%%%%%%%%%%%%%%%%%%%%%%%%%%%%%%%%%%%%%%%%%%%%

%\end{document}
